
\textbf{1. Definir as dimensões das grandezas}

Primeiro, identificamos as dimensões fundamentais de cada variável envolvida no sistema (Massa [M], Comprimento [L] e Tempo [T]):
\begin{itemize}
    \item Tensão (T): É uma força, logo $[T] = [M][L][T]^{-2}$
    \item Comprimento (L): $[l] = [L]$
    \item Massa (m): $[m] = [M]$
    \item Velocidade (v): $[v] = [L][T]^{-1}$
\end{itemize}

\textbf{2. Estabelecer a equação dimensional}

Assumimos que a tensão $T$ depende de $m$, $v$ e $l$ elevadas a potências desconhecidas $a$, $b$ e $c$:
\[
    [T] = k \cdot m^a \cdot v^b \cdot l^c
\]
Substituindo pelas dimensões:
\[
    [M]^1[L]^1[T]^{-2} = ([M])^a ([L][T]^{-1})^b ([L])^c
\]
\[
    [M]^1[L]^1[T]^{-2} = [M]^a [L]^{b+c} [T]^{-b}
\]

\textbf{3. Resolver o sistema de equações}

Igualando os expoentes de cada dimensão fundamental:
\begin{align*}
    &\text{Para } [M]: & 1 &= a \Rightarrow a = 1 \\
    &\text{Para } [T]: & -2 &= -b \Rightarrow b = 2 \\
    &\text{Para } [L]: & 1 &= b + c \Rightarrow 1 = 2 + c \Rightarrow c = -1
\end{align*}

\textbf{4. Formar o grupo adimensional}

\[
    [T] = k \cdot m^1 v^2 l^{-1}
\]
\[
    T = k \frac{mv^2}{l}
\]

Isolando a constante $k$ para encontrar o grupo adimensional pedido:
\[
    \frac{Tl}{mv^2} = k
\]

Como $k$ é uma constante numérica (adimensional), provamos que a relação entre as variáveis é governada por esse grupo.
